\documentclass[11pt]{article}
\usepackage[margin=1in]{geometry}
\usepackage{xcolor}
\usepackage{tcolorbox}
\usepackage{hyperref}
\usepackage{enumitem}
\setlist{noitemsep,topsep=0pt}
\usepackage{booktabs}
\usepackage{array}
\usepackage{amsmath}

% Define OSU orange and AI blue colors
\definecolor{osaorange}{RGB}{220,115,38}
\definecolor{aiblue}{RGB}{30,144,255}

% Configure hyperref colors
\hypersetup{
    colorlinks=true,
    linkcolor=blue,
    urlcolor=blue,
    citecolor=blue
}

\title{}
\author{}
\date{}

\begin{document}

% Orange header box with black outline
\begin{tcolorbox}[colback=osaorange,colframe=black,boxrule=2pt,arc=0pt,left=6pt,right=6pt,top=10pt,bottom=10pt]
\centering
\textcolor{white}{\textbf{\Large Week 3 In-Class Worksheet}}\\
\textcolor{white}{\textbf{\large Probability Distributions and Sampling Distributions}}
\end{tcolorbox}

\vspace{8pt}

\noindent\textbf{Name:} \rule{8cm}{0.4pt} \hfill \textbf{Date:} \rule{4cm}{0.4pt}\\[0.5ex]
\textbf{Course:} H524 Introduction to Biostatistics\\
\textbf{Topic:} Binomial, Normal, Sampling Distributions, and Confidence Intervals

\vspace{8pt}

\section*{Instructions}
Work with a partner to solve the following problems. Show all your work. You may use R to verify your answers, but demonstrate your understanding of the concepts by showing calculations.

\vspace{12pt}

\section{Problem 1: Binomial Distribution}

A new antibiotic has an 80\% success rate in treating a particular bacterial infection. A clinic treats 15 patients with this antibiotic.

\vspace{6pt}

\textbf{Part A:} What is the probability that exactly 12 patients are successfully treated?

\vspace{0.3cm}

\textbf{Given:} $n = 15$, $p = 0.80$, find $\Pr(X = 12)$

\vspace{0.3cm}

\textbf{Formula:} $\Pr(X = k) = \binom{n}{k} p^k (1-p)^{n-k}$

\vspace{3.5cm}

\textbf{Part B:} What is the probability that at least 13 patients are successfully treated?

\vspace{0.3cm}

Hint: $\Pr(X \geq 13) = \Pr(X=13) + \Pr(X=14) + \Pr(X=15)$

\vspace{3.5cm}

\textbf{Part C:} Calculate the expected number of successful treatments and the standard deviation.

\vspace{0.3cm}

Expected value: $\mu = np$

\vspace{1cm}

Standard deviation: $\sigma = \sqrt{np(1-p)}$

\vspace{1.5cm}

\textbf{Part D:} Use R to verify your answers. Write the R code you would use:

\vspace{2.5cm}

\newpage

\section{Problem 2: Normal Distribution}

Adult systolic blood pressure is normally distributed with mean 125 mmHg and standard deviation 18 mmHg.

\vspace{6pt}

\textbf{Part A:} What proportion of adults have blood pressure above 140~mmHg (hypertension)?

\vspace{0.3cm}

\textbf{Step 1:} Standardize using $Z = \frac{X - \mu}{\sigma}$

\vspace{2cm}

\textbf{Step 2:} Find $\Pr(Z > z)$ using the standard normal table or R

\vspace{2cm}

\textbf{Part B:} What proportion have blood pressure between 110 and 140 mmHg?

\vspace{0.3cm}

Hint: $\Pr(110 < X < 140) = \Pr(Z < z_2) - \Pr(Z < z_1)$

\vspace{3.5cm}

\textbf{Part C:} Find the blood pressure value that represents the 90th percentile.

\vspace{0.3cm}

\textbf{Step 1:} Find $z$ such that $\Pr(Z < z) = 0.90$

\vspace{1.5cm}

\textbf{Step 2:} Convert back to original scale: $X = \mu + z\sigma$

\vspace{2cm}

\textbf{Part D:} Write the R code to solve parts A, B, and C:

\vspace{2.5cm}

\newpage

\section{Problem 3: Sampling Distribution and Central Limit Theorem}

A population of cholesterol levels has mean $\mu = 200$ mg/dL and standard deviation $\sigma = 40$ mg/dL. A sample of $n = 64$ individuals is randomly selected.

\vspace{6pt}

\textbf{Part A:} What is the distribution of the sample mean $\bar{X}$?

\vspace{0.3cm}

By the Central Limit Theorem: $\bar{X} \sim N(\mu_{\bar{X}}, \sigma_{\bar{X}}^2)$

\vspace{0.3cm}

Mean of $\bar{X}$: $\mu_{\bar{X}} = $

\vspace{1cm}

Standard error of $\bar{X}$: $\sigma_{\bar{X}} = \frac{\sigma}{\sqrt{n}} = $

\vspace{1.5cm}

\textbf{Part B:} What is the probability that the sample mean is greater than 210 mg/dL?

\vspace{0.3cm}

\textbf{Step 1:} Standardize: $Z = \frac{\bar{X} - \mu_{\bar{X}}}{\sigma_{\bar{X}}}$

\vspace{2cm}

\textbf{Step 2:} Find probability

\vspace{2cm}

\textbf{Part C:} What is the probability that the sample mean is between 195 and 205 mg/dL?

\vspace{3.5cm}

\textbf{Part D:} How would your answers change if $n = 16$ instead of $n = 64$? Would that probability increase or decrease? Explain why.

\vspace{3cm}

\newpage

\section{Problem 4: Confidence Intervals}

A clinical trial measures the reduction in pain scores for 20 patients receiving a new medication. The sample mean reduction is $\bar{x} = 4.5$ points, with sample standard deviation $s = 1.8$ points.

\vspace{6pt}

\textbf{Part A:} Calculate a 95\% confidence interval for the true mean pain reduction.

\vspace{0.3cm}

\textbf{Step 1:} Determine what critical value to use. Should you use $z$ or $t$? Why?

\vspace{2cm}

\textbf{Step 2:} Find the critical value.

\vspace{0.3cm}

Degrees of freedom = $n - 1 = $

\vspace{0.5cm}

Critical value: $t_{0.975, df} = $ \hspace{1cm} (use~R: \texttt{qt(0.975,~df)})

\vspace{1.5cm}

\textbf{Step 3:} Calculate the standard error: $\text{SE} = \frac{s}{\sqrt{n}}$

\vspace{2cm}

\textbf{Step 4:} Calculate the margin of error: $\text{ME} = t_{critical} \times \text{SE}$

\vspace{2cm}

\textbf{Step 5:} Construct the confidence interval: $\bar{x} \pm \text{ME}$

\vspace{2.5cm}

\textbf{Part B:} Interpret the confidence interval in context. Write a complete sentence explaining what this interval means.

\vspace{3cm}

\textbf{Part C:} If you wanted a 99\% confidence interval instead, would it be wider or narrower than the 95\% interval? Explain.

\vspace{2.5cm}

\newpage

\section{Problem 5: Integration Problem}

A pharmaceutical company is testing a new cholesterol-lowering drug. In the clinical trial:
\begin{itemize}
\item 100 patients receive the drug
\item Cholesterol reduction follows a normal distribution
\item Sample mean reduction: $\bar{x} = 28$ mg/dL
\item Sample standard deviation: $s = 12$ mg/dL
\end{itemize}

\vspace{6pt}

\textbf{Part A:} Calculate a 95\% confidence interval for the true mean cholesterol reduction.

\vspace{4cm}

\textbf{Part B:} Based on your confidence interval, is there strong evidence that the drug reduces cholesterol by more than 25 mg/dL on average? Explain your reasoning.

\vspace{3cm}

\textbf{Part C:} If individual cholesterol reductions are normally distributed with the sample statistics above, approximately what proportion of patients will experience a reduction of at least 40 mg/dL?

\vspace{0.3cm}

Hint: Use the sample statistics as estimates of population parameters.

\vspace{4cm}

\textbf{Part D:} Explain the difference between the questions in Parts B and C. Why do we use different approaches (confidence interval vs.\ normal probability)?

\vspace{3.5cm}

\vspace{1cm}

\section*{Reflection Questions}

\textbf{1.} What is the key difference between a standard deviation and a standard error?

\vspace{2cm}

\textbf{2.} When should you use a $t$-distribution instead of a normal distribution for confidence intervals?

\vspace{2cm}

\textbf{3.} How does sample size affect the width of a confidence interval?

\vspace{2cm}

\end{document}
